% OLD TEXT FROM HERE ON
% This chapter contains text from the original submission.
% It will be revised as part of the salvage plan.

\section{Hacking demand: varieties of labor market slack} \label{sec:labor}

Based on the preliminary findings of the preceding section we now move on to consider the impact of a second major \textit{axiomatic variation} \citep{Kapeller.2013} proposed in the papers of Baqaee and Farhi, especially in \citet{Baqaee.2021,BF2022}, which relates to the introduction of underutilized factor markets, that is, slack on labor markets. This variation allows for integrating the notion of (involuntary) unemployment, which affects (capacity) constraints on production factors. These constraints play, as we have seen, a key role for determining the overall magnitude of the results presented so far. Similar to other classic examples of axiomatic variation in economics (like the "market for lemons" in the early 1970s), the authors provide an overall narrative or story (akin to the reference to the market for used cars) to create legitimacy for looking at ideas running counter to established convictions and conventions (like asymmetric information and adverse selection; see \cite{Akerlof.1970}). In this application the narrative legitimizing the consideration of such labor market slack is given by the context of Corona-related lockdowns that partially required people to abstain from work due to social distancing regulations.\footnote{The attentive reader will identify a second source of unemployment mentioned in \citet{BF2022}, which follows from the assumption of minimum wages leading to market inefficiencies. However, this form of unemployment does not relax capacity constraints, which is why we disregard this aspect in our analysis.}

The introduction of such a labor slack in this simple model setup implies -- also because the model features only one production factor -- a fundamental relaxation of the crucial capacity constraint, i.e., the conceptual source of scarcity in the model. Against this backdrop, it does not come as a great surprise that the general magnitude of observed outcomes in terms of GDP crucially depend on the amount of labor slack allocated in the first place, while the sector-specific reactions in terms of prices and strongly depend on the sectoral allocation of labor slack. 

Conceptually, it is important to note that the introduction of such a labor slack by itself is to be conceived and interpreted as the advent of an exogenous shock: in model terms the economy is at some variant of full employment already in the pre-shock state, whereas the assumption of labor slack kicks in \textit{jointly} with other assumed shocks and effectively implies the arrival or emergence of additional labor inputs, which are, in turn, absorbed by efficient markets to produce additional outputs. Hence, even in the absence of other shocks -- like the demand shock in our application -- the introduction of labor slack would imply the search for a equilibrium as it effectively represents a spontaneous expansion of the labor force.

In what follows, we employ two different rationales to be potentially applied to the case at hand: first, we consider a theoretical rationale in section \ref{subsec:capconstraint1} that starts from assuming that labor markets can mobilize workers to match the necessary expansion in a case of a demand shock. For doing so, we take the labor market expansion as predicted by the simple Leontief model as a benchmark to guarantee an adequate degree of slack to be imposed in a CGE-context. Second we consider a naive empirical calibration approach that employs (some fraction of) current unemployment as a proxy for existing slack on labor markets in section \ref{subsec:capconstraint2}. In both applications we specify the relevant expansion of the labor force on a sectoral level (i.e., as a vector), which either mimics the impact of the shock as predicted by the Leontief model (as in section \ref{subsec:capconstraint1}) or increases the labor force allocated to all sectors by some constant growth rate (as in section \ref{subsec:capconstraint2}).

\subsection{Abstracting from capacity constraints I: a demand-driven CGE model?} \label{subsec:capconstraint1}

To explore the practical implications of labor slack, we use a CES-based CGE model as outlined in section \ref{subsec:introducingCES} that calibrates labor slack by taking the labor market expansion associated with a traditional Leontief approach as a benchmark This procedure allocates labor in accordance with projected sectoral requirements and leads to an aggregate employment growth of 1.7\%. 

When inspecting the model results we observe at first glance that the theoretical intuition outlined above -- namely that capacity constraints prove crucial for output adjustment -- is immediately confirmed. Indeed, we find that by approaching the problem this way the results of the CGE model for real GDP growth are much more in line with the traditional Leontief outcome, even surpassing the latter slightly. While the impact of the exact choices of elasticities remains negligible for the larger part of possible parameter values, we find that the introduction of labor slack -- or more generally slack on factor markets -- indeed provides a preliminary vantage point for constructing meaningful conceptual connections between both modeling traditions. Crucially, the general magnitude of overall GDP is seemingly governed by the amount of additional labor allocated, where the former estimate is close to, but generally (slightly) below, the size of the imposed labor slack.

\begin{figure}
    \centering
    \includegraphics[width=1\linewidth]{pictures/impulse_ls.png}
    \caption{Effect of different elasticities on GDP with labor slack (Leontief calibration). Note that, in the visual display there is a strong overlay between all scenarios associated using a CES-function and the Cobb-Douglas model.}
    \label{fig:elasticity-gradient2}
\end{figure}

In contrast to this significant overlap in terms of aggregate outcomes, a key difference between a classic Leontief input-output model and the CGE-version, that mimics the overall growth and employment effects of the former, is that the latter also allows for greater heterogeneity of sectoral outcomes in terms of both, quantities and prices. Overall we find that -- similar to the Leontief-case -- all sectors in the economy will expand their production, whereas price changes are somewhat idiosyncratic. Specifically, we find that quantity adjustments across sectors are largely governed by the additional labor allocated to those sectors. Relative prices, however, do not correlate strongly with allocated labor, but rather indicate the presence of trade-offs. For the shocked sectors a stronger quantity expansion tends to co-occur with a greater price increase (with one exceptional outlier, see the right panel in Figure \ref{fig:panel_ls}). In contrast, the inverse pattern emerges for sectors not affected from the shock (see the right panel in Figure \ref{fig:panel_ls}). In this context, increasing prices again imply crowding-out effects for shocked sectors in terms of quantity produced, while those shocked sectors that find prices decreasing are also subject to a corresponding crowding-in effect in terms of quantity to match the desired output adjustment.

These differences in post-shock relative equilibrium prices and quantities across sectors are thereby driven by the distinct use of intermediate goods: for shocked sectors the startling exception is glass, which combines a strong quantity expansion with a price decrease. This outcome is due to its somewhat distinct structure of intermediates as compared to the other shocked sectors, which lets it profit from cheap relative prices. Similarly, the sector with the strongest price increase -- ceramics -- needs inputs in the form of primary goods, that are seemingly originating from sectors more difficult to expand, which in turn leads to over-proportional price increases.

Overall the origin of the observed changes in relative prices have to be interpreted against the backdrop that the additional labor allocated on a sectoral level is taken from an extrapolation based on a Leontief-projection. This 'enforced' sectoral allocation of labor governs quantities that seemingly deviate from a general equilibrium allocation and translate into an oversupply of labor in some sectors (leading to lower prices) and a shortage in others (which raises prices). This misalignment leads to small losses in terms of efficiency, which is why growth estimates fall slightly behind the allocate amount of labor slack.

However, we nonetheless observe that in this scenario, where labor market slack is aligned with the exogenous shock, volatility in relative prices is much smaller and more symmetric than in all other scenarios analyzed in this paper. In other words, the release of additional labor supply incorporated by the labor slack assumptions effectively prevents any substantial overheating of the economy and has rather symmetric effects on sectorial efficiency. 

\begin{figure}
    \centering
    \includegraphics[width=1\linewidth]{pictures/panel_ls.png}
    \caption{Changes in prices and consumption for shocked sectors (left panel) and quantities and prices for all sectors (right panel). Outputs are based on a model run using the standard specification from \citet{BF2019}, i.e., $(\varepsilon,\theta,\sigma) = (0.001, 0.5, 0.9)$, with added labor slack from the Leontief approximation. Error bars indicate the range of outcomes that can be achieved under for extreme parameter values (specifically, $(\varepsilon,\theta,\sigma) = (0.001, 0.1, 0.6)$ and $(\varepsilon,\theta,\sigma) = (0.99, 0.99, 0.99)$).}
    \label{fig:panel_ls}
\end{figure}

When comparing the expected adjustments in gross output implied by the Leontief projection and the CGE model with labor slack across all sectors as shown in Figure \ref{fig:sectoral-changes_ls} we observe that the latter version more strongly expands the production of intermediate goods \textit{on average}, but not consistently across all sectors. The reason is for this seemingly unsystematic outcome is that the CGE model identifies primarily those sectors as constrained in their expansion that are either associated with the provision of primary goods, like oil and gas and metallic ores, or services (large remainder). In other words, the CGE shows two major qualitative differences in its projection: first, it substitutes away from some primary sectors and compensates this with a significant expansion of intermediary goods in related sectors (especially building construction work, civil engineering work, cars and trucks, refined metals of various sorts as well as machinery). Second, it implicitly assumes a much smaller or qualitatively different reaction of consumers to receiving additional income, which leads to comparatively smaller expenditures for large parts of the service sectors.

\begin{figure}
    \centering
    \includegraphics[width=1\linewidth]{pictures/diff_lambda_imp_ls.png}
    \caption{Change in outputs, based on a model run using the standard specification from \citet{BF2019}, i.e., $(\varepsilon,\theta,\sigma) = (0.001, 0.5, 0.9)$ with added labor slack from the Leontief approximation. Shocked sectors remain unshaded.}
    \label{fig:sectoral-changes_ls}
\end{figure}

\subsection{Abstracting from capacity constraints II: adopting an empirical calibration?}\label{subsec:capconstraint2}

As already indicated, in this second application we again assume that the shock to the economy has two components: the imposed shift in final demand as well as the advent of additional labor, that was previously unavailable, thereby increasing the total available labor force. However, in this application the allocation of additional labor mimics the existing sectoral allocation, i.e., the additional labor is transferred uniformly to the sectors, so the new labor supply is simply given by $\mathbf l_\text{new} = (1 - \mu)^{-1} \mathbf l$. While the choice of $\mu$ is flexible in our baseline application $\mu$ is equal to the rate of unemployment.\footnote{Unemployment in Germany in 2019 according stood at 3.2 \% according to the ILO definition of unemployment, which corresponds to roughly 60\% of registered unemployment in Germany as provided by Federal employment statistics  (\cite{LabourMarketInformation}).}

Hence, this second empirical application is based on a different -- and conceptually somewhat more independent -- reaction of labor markets, where the sectoral expansion of workers is not directly tied to (estimated) requirements of the respective demand shock. At the same time the assumed reaction is more sizeable as total employment expands by 3.1\%, whereas in the former application the total expansion of the labor force was only at 1.7\%.

\begin{figure}
    \centering
    \includegraphics[width=1\linewidth]{pictures/impulse_ls_empirical.png}
    \caption{Effect of different elasticities on GDP with labor slack (empirical labor slack)}
    \label{fig:elasticity-gradient3}
\end{figure}


Again we observe that the overall magnitude of growth estimated by the CGE model is largely governed by the allocated labor slack as is illustrated by Figure \ref{fig:panel_ls_emp}. However, the misalignment between exogenous changes in final demand and available factors is reflected by the fact that GDP results are more sensitive to parameter choices on elasticities as compared to the application in section \ref{subsec:capconstraint1}. Also this misalignment reappears on a sectoral level, so that the corresponding results reproduce the patterns already found in section \ref{sec:ces} .

Nonetheless, on a general level this example again illustrates how the introduction of labor slack can lead to more optimistic results in terms of GDP growth. The key requirement is the willingness to assume that such an expansion of the labor force is indeed a plausible reaction -- either due to labor market considerations or as a general response to increased effective demand. While this feature is not truly surprising it shows how modifying dominant convictions on the role and modeling of capacity constraints indeed opens up a shared realm, in which the simplistic juxtaposition of an 'expansionary' Leontief-approach and a 'contractionary' CGE-perspective becomes subject to a more nuanced perspective that allows for pragmatically interpreting labor slack as a moderating variable to explore some shared ground between these two competing visions of the economy.


\begin{figure}
    \centering
    \includegraphics[width=1\linewidth]{pictures/panel_ls_empirical.png}
    \caption{Changes in prices and final consumption for shocked sectors (left panel) and quantities and prices for all sectors (right panel). Outputs are based on a model run using the standard specification  from \citet{BF2019}, i.e. $(\varepsilon,\theta,\sigma) = (0.001, 0.5, 0.9)$, with added labor slack (3.1\% per sector). Error bars indicate the range of outcomes that can be achieved under for extreme parameter values (specifically, $(\varepsilon,\theta,\sigma) = (0.001, 0.1, 0.6)$ and $(\varepsilon,\theta,\sigma) = (0.99, 0.99, 0.99)$).}
    \label{fig:panel_ls_emp}
\end{figure}

A deeper analysis of the results of this application shows that the underlying assumption on how exactly additional labor is allocated matters greatly for the sectoral outcomes predicted by the CGE-approach. In subsection \ref{subsec:capconstraint1} assigning large segments of the additional labor force explicitly to shocked sectors led to a complex reconfiguration of the economy that shifted intersectoral relationships substantially and in heterogeneous ways, but managed to accompany the shock in final demand quite well (as is evidenced by the moderate price increase in shocked sectors). In contrast, in this application the economy now expands more homogeneously across sectors (see Figure \ref{fig:sectoral-changes_ls_emp} for an overview), but falls short in terms of final outputs produced in shocked sectors, which also see relative prices increase strongly. These sectoral bottlenecks in turn create an overall pattern of quantity and price effects similar to the original application without labor slack (see especially the right panel in Figure \ref{fig:panel_ls_emp} in comparison with the right panel of Figure \ref{fig:comparison-sec4} in section \ref{subsec:CGE_elasticity}) that is, however, coined by an overall expansion of production, which is in large parts governed by the amount of labor slack allocated. 

The much more symmetric expansion in output across all sectors is very close to the result observed in section \ref{subsec:CGE_elasticity} (e.g. in Figure \ref{fig:panel_ls_emp}), where no labor slack has been imposed at all. Hence, we observe that the allocating additional labor supply in strict proportion to current employment across sector will uniformly increase produced quantities, but will, otherwise, only show a minimal impact on prices. A shock on final demand, on the other hand, might easily have a stronger impact on prices than quantities even in the case of labor slack (see the left panel in Figure \ref{subsec:CGE_elasticity}), if some sectorial misalignment remains.

This latter observation also translates into a sectorial view as indicated by Figure \ref{fig:sectoral-changes_ls_emp}. Here we observe that the expansion of total production falls behind the predictions of the Leontief model for the shocked sectors, while most other sectors show a significantly greater expansion than before. A few exceptions are again found in sectors, which focus on the provision of primary goods, like ores, wood and refined petroleum and which are substituted by alternative intermediaries in the new equilibrium.

\begin{figure}
    \centering
    \includegraphics[width=1\linewidth]{pictures/diff_lambda_imp_ls_empirical.png}
    \caption{Change in outputs, based on a model run using the standard specification from \citet{BF2019}, i.e. $(\varepsilon,\theta,\sigma) = (0.001, 0.5, 0.9)$, with added labor slack (3.1\% per sector). Shocked sectors remain unshaded.}
    \label{fig:sectoral-changes_ls_emp}
\end{figure}


